% ==============================================================================
% FRONTMATTER: ABSTRACT.TEX
% Tóm tắt báo cáo học phần và từ khóa nghiên cứu (Keywords).
% ==============================================================================

\chapter*{TÓM TẮT BÁO CÁO}
\addcontentsline{toc}{chapter}{TÓM TẮT BÁO CÁO}

Báo cáo học phần \course\ với đề tài \textbf{``\reporttitle''} trình bày quá trình nghiên cứu, thiết kế và phát triển hoàn chỉnh một nền tảng phần mềm hỗ trợ rèn luyện thể hình và theo dõi chế độ dinh dưỡng ứng dụng các kỹ thuật trí tuệ nhân tạo hiện đại. Trong bối cảnh nhu cầu tập luyện cá nhân hóa ngày càng tăng cao nhưng chi phí thuê huấn luyện viên riêng còn đắt đỏ, đồng thời người tự tập thường thiếu kiến thức cơ sinh học dẫn đến nguy cơ chấn thương hoặc tập luyện quá tải, đề tài đã đề xuất giải pháp kết hợp giữa năng lực phân tích dữ liệu tức thời của máy tính và quy trình giám sát chuyên môn chặt chẽ từ huấn luyện viên thể hình thực tế.

Hệ thống được xây dựng theo kiến trúc phân tầng hiện đại gồm giao diện người dùng React 19, máy chủ ứng dụng Node.js Express 5, cơ sở dữ liệu quan hệ PostgreSQL và hạ tầng suy luận trí tuệ nhân tạo lai giữa máy chủ Ollama cục bộ và Google Gemini API. Trọng tâm của đề tài bao gồm ba đóng góp kỹ thuật chính: thứ nhất, phát triển thuật toán tự động tính toán nhu cầu năng lượng nền tảng và phân bổ dưỡng chất đa lượng theo công thức Mifflin-St Jeor kết hợp cơ chế sinh giáo án bốn pha an toàn chấn thương; thứ hai, tích hợp kỹ thuật RAG nhằm loại trừ sai lệch thông tin khi tư vấn kỹ thuật vận động; và thứ ba, ứng dụng kỹ thuật LoRA, QLoRA để tinh chỉnh mô hình thị giác máy tính nhận diện miền ảnh và bóc tách thành phần dinh dưỡng của các món ăn Việt Nam. Đặc biệt, hệ thống thiết lập nguyên tắc phân quyền cốt lõi là quyết định của huấn luyện viên luôn có quyền ưu tiên cao hơn trí tuệ nhân tạo, đảm bảo an toàn tối đa cho học viên.

Kết quả kiểm thử thực nghiệm trên các kịch bản thực tế chứng minh hệ thống vận hành ổn định, thời gian xử lý cơ sở dữ liệu duy trì dưới 150 ms, tốc độ suy luận mô hình ngôn ngữ và thị giác đạt yêu cầu tương tác thời gian thực, đồng thời tỷ lệ nhận diện món ăn Việt đạt trên 91\% và tỷ lệ loại bỏ bài tập chống chỉ định chấn thương đạt 100\%. Đề tài đã hoàn thành toàn diện các mục tiêu đề ra và mở ra tiềm năng ứng dụng thiết thực trong việc số hóa quy trình vận hành tại các trung tâm thể hình hiện đại.

\vspace{1.0cm}
\noindent\textbf{Từ khóa:} Trí tuệ nhân tạo, Huấn luyện viên ảo, Mô hình ngôn ngữ lớn, Thị giác máy tính, Tinh chỉnh LoRA, RAG, Quản lý lịch tập thể hình và dinh dưỡng.

\clearpage
